\section{Interferometric products (INT / COHE / UNW \& GINT / GCOHE / GUNW)}

\subsection{Product Description}
The Interferogram (\acs{int}) product represents the wrapped phase difference (in radians) between two co-registered SAR acquisitions and is used to analyze surface deformation and coherence. It is generated from co-registered \acs{slc}s or \acs{cslc}s, with corrections applied for co-registration errors (e.g., Enhanced Spectral Diversity) and, when necessary, ionospheric effects (e.g., via split-spectrum technique).

The \acs{int} product includes:
\begin{itemize}
    \item Wrapped phase information
    \item Normalized interferometric coherence (\acs{cohe})
    \item Auxiliary layers such as DEM, baseline lookup tables, and optional phase correction maps
\end{itemize}

The product is available in both:
\begin{itemize}
    \item \textit{Radar geometry} (burst-wise or swath-merged)
    \item \textit{Ground-projected format} (\acs{gint} and \acs{gcohe}) in projected coordinates
\end{itemize}

The resolution may vary depending on the application, from full-resolution (e.g., urban ground deformation) to coarser, multilooked resolutions for regional-scale analyses.

The Unwrapped Interferogram (\acs{unw}) product contains the unwrapped phase (in radians), where the 2$\pi$ ambiguities present in the wrapped interferogram have been resolved. This product enables direct estimation of surface displacement, provided that the phase unwrapping process is successful and robust.

The \acs{unw} product includes:
\begin{itemize}
    \item Unwrapped phase
    \item Coherence map (\acs{cohe})
    \item Geometry masks and connected components
    \item Auxiliary layers, baseline lookup tables, and optional phase correction maps
\end{itemize}

\acs{unw} can be delivered in:
\begin{itemize}
    \item \textit{Radar geometry} (as \acs{unw})
    \item \textit{Ground-projected format} (\acs{gunw} and \acs{gcohe}) in projected coordinates
\end{itemize}

Because phase unwrapping is a known ill-posed problem and no fully reliable algorithm exists that universally guarantees high-quality results, the dissemination of \acs{unw} products must be carefully assessed. Their generation is typically limited to areas where sufficient coherence and signal stability allow for accurate unwrapping.

The interferometric products are derived from pairs of SAR acquisitions. These products are planned to be generated specifically from pairs of consecutive or time-close acquisitions.

\acs{int}, \acs{cohe}, and \acs{unw} will be generated using \acs{cslc}s and are defined on a common reference radar grid shared by both the Sentinel-1 and ROSE-L missions. The output radar grid of these products will likewise be aligned across both missions, ensuring interoperability and joint usability of products generated from data acquired by either satellite.

\subsubsection{Input and Auxiliary Data}
The following input and auxiliary data are considered for product generation:

\begin{longtblr}[
    label = tab:in_aux_int_products,
    entry = none,
    caption = {Interferometric products' input and auxiliary data},
    ]{
    width = \linewidth,
    colspec = {Q[l,20]Q[l,60]Q[c,20]},
    hlines,
    hline{1-Z} = {-}{0.05em},
    }
    \textbf{Name}   & \textbf{Description}   & \textbf{Mandatory (M) /       Optional (O)} \\
    Reference \acs{cslc} & Reference L1 Co-registered Single Look
    Complex Image and annotations & M \\
    Secondary \acs{cslc} & Secondary L1 Co-registered Single Look
    Complex Image and annotations & M \\
    \acs{dem} & External DEM (e.g., Copernicus DEM at different resolutions, with 30 m as the default worldwide). & M \\
    Precise Orbits & \acs{cslc} precise orbits. & M \\
    \acs{etad} & \acs{etad}-like products for atmospheric delay and timing corrections. & O \\
    \acs{tec} & External \acs{tec} data for delay correction & O \\
\end{longtblr}

\subsubsection{Main Variables}
The following tables list the main variables associated with the interferometric product.

\begin{longtblr}[
    label = tab:var_int_products,
    entry = none,
    caption = {Interferometric product main variables},
    ]{
        width = \linewidth,
        colspec = {Q[l,25]Q[l,52]Q[l,20]Q[l,13]},
        hlines,
        hline{1-Z} = {-}{0.05em},
    }
    \textbf{Name} & \textbf{Description} & \textbf{Dimensions} & \textbf{Data Type} \\
    interferogram & Wrapped interferometric phase between Secondary and Reference \acs{cslc}s, corrected for the topographic contribution and, optionally, for ionospheric and atmospheric effects. Present in \acs{int} and \acs{unw} and their projected version & polarization and 
    
    - range\_time, azimuth\_time for radar geometry products
    
    - x, y for projected products & float32 \\
    unwrapped 
    
    \_interferogram & Unwrapped interferometric phase between the Secondary and Reference \acs{cslc}s, corrected for the topographic contribution and, optionally, for ionospheric and atmospheric effects. Present in \acs{unw} and its projected version & polarization and 
    
    - range\_time, azimuth\_time for radar geometry products
    
    - x, y for projected products & float32 \\
    coherence & The coherence of between the reference and secondary \acs{cslc}s forming the interferometric pair & polarization and 
    
    - range\_time, azimuth\_time for radar geometry products
    
    - x, y for projected products & float32 \\
\end{longtblr}

The table below describes the dimensions/coordinates associated with the interferometric product main variables.

\begin{longtblr}[
    label = tab:var_int_coordinates,
    entry = none,
    caption = {Interferometric product main variables coordinates/dimensions},
    ]{
        width = \linewidth,
        colspec = {Q[l,20]Q[l,62]Q[l,18]},
        hlines,
        hline{1-Z} = {-}{0.05em},
    }
    \textbf{Name} & \textbf{Description} & \textbf{Data Type} \\
    polarization & Acquired polarizaions & string \\
    azimuth\_time & Time along satellite track direction & datetime64 \\
    range\_time & Time across satellite track direction & float64 \\
    x & UTM/UPS x (East) coordinates [m] & float64 \\
    y & UTM/UPS y (North) coordinates [m] & float64 \\
\end{longtblr}

\subsubsection{Main Metadata}
The following tables list the main metadata and per-pixel metadata associated with the interferometric products.

\begin{longtblr}[
    label = tab:meta_int_products,
    entry = none,
    caption = {Interferometric products' main metadata and per-pixel metadata},
    ]{
    width = \linewidth,
    colspec = {Q[l,17]Q[l,38]Q[l,20]Q[l,13]Q[c,12]},
    hlines,
    hline{1-Z} = {-}{0.05em},
    }
    \textbf{Name} & \textbf{Description} & \textbf{Dimensions} & \textbf{Data Type} & \textbf{Per-pixel Metadata}\\
    ionospheric
    
    \_phase & The ionospheric phase subtracted from the phase, if performed &  
    - range\_time, azimuth\_time for radar geometry products
    
    - x, y for projected products &  float32 & \ding{51} \\
    tropospheric
    
    \_phase & The tropospheric phase subtracted from the phase, if performed &  
    - range\_time, azimuth\_time for radar geometry products
    
    - x, y for projected products &  float32 & \ding{51} \\
    topographic
    
    \_phase & Topographic phase estimated by using DEM & - range\_time, azimuth\_time for radar geometry products
    
    - x, y for projected products & float32 & \ding{51} \\
    reference\_cslc & Link to input reference \acs{cslc} & - & string & \\
    secondary\_cslc & Link to input secondary \acs{cslc} & - & string & \\
    reference\_orbit & Link to reference orbit file & - & string & \\
    reference\_dem & Link to reference DEM & - & string & \\
    etad\_file & ETAD file used for each image belonging to the CSLC stack. This is optional: it is included only if used; otherwise, it is null & - & string & \\ 
    etad\_azcorr & Extended Timing Annotation Data: azimuth time correction resampled onto the common grid of the reference geometry. This is optional: it is included only if used; otherwise, it is null & 
    - azimuth\_time, slant\_range\_time  for product in radar geometry 
    
    - x, y for projected products & float64 & \ding{51} \\
    etad\_rgcorr & Extended Timing Annotation Data: slant range time correction resampled onto the common grid of the reference geometry. This is optional: it is included only if used; otherwise, it is null & - azimuth\_time, slant\_range\_time  for product in radar geometry 
    
    - x, y for projected products & float64 & \ding{51} \\
    tec & Link to external \acs{tec} data, if used & - & string & \\ 
    look\_vector & 3-D components radar unit look vector, in an Earth-Centred Earth-Fixed (\acs{ecef}) coordinate system. It consists of unit vectors from antenna to surface pixel. Only for reference geometry. &  (x, y) or (azimuth\_time, range\_time) & float64 & \ding{51} \\
    % los\_vector & Director cosines of the line of sight & - azimuth\_time, slant\_range\_time  for product in radar geometry 
    
    % - x, y for projected products & float32 & \ding{51} \\
    mask & Mask containing, for pixel, the following information:
    
     - 0: no distortions
    
     - 1: layover
    
     - 2: shadow
    
     - 3: undefined & slant\_range\_time  for product in radar geometry 
    
    - x, y for projected products & byte & \ding{51} \\
     
    ml\_fact & Applied range-azimuth multilook factor & slant\_range\_time  for product in radar geometry 
    
    - x, y for projected products & float32 & \ding{51} \\
    x\_sample & Pixel spacing in azimuth or North direction for radar geometry and projected products, respectively & - & float32 & \\
    y\_sample & Pixel spacing in slant range or East direction for radar geometry and projected products, respectively & - & float32 & \\
    \_geometric
    
    \_accuracy
    & Expected geometric accuracy, including eastern and northern geometric accuracy bias and standard deviation. Only for projected products. The measures will be performed on a validation dataset and periodically updated. & - & dictionary & \\ 
    east\_geom
    
    \_bias & Expected eastern geometric bias.  Only for projected products & - & float32 & \\
    east\_geom
    
    \_std & Expected eastern geometric standard deviation.  Only for projected products & - & float32 & \\
    north\_geom
    
    \_bias & Expected northern geometric bias.  Only for projected products & - & float32 & \\
    north\_geom
    
    \_std & Measured northern geometric standard deviation.  Only for projected products & - & float32 & \\
    geometric
    
    \_accuracy
    
    \_radial
    
    \_rmse	& Expected radial RMSE of geometric accuracy. Only for projected products & - & float32 & \\
    
\end{longtblr}
    
The table below describes the dimensions/coordinates associated with the interferometric products' main metadata.

\begin{longtblr}[
    label = tab:var_int_metadata_coordinates,
    entry = none,
    caption = {Interferometric products' main metadata coordinates/dimensions},
    ]{
    width = \linewidth,
    colspec = {Q[l,20]Q[l,62]Q[l,18]},
    hlines,
    hline{1-Z} = {-}{0.05em},
    }
    \textbf{Name} & \textbf{Description} & \textbf{Data Type} \\
    polarization & Acquired polarizations & string \\
    azimuth\_time & Time along satellite track direction & datetime64 \\
    range\_time & Time across satellite track direction & float64 \\
    x & UTM/UPS x (East) coordinates [m] & float64 \\
    y & UTM/UPS y (North) coordinates [m] & float64 \\
\end{longtblr}

\subsubsection{Reference System Tiling and Slicing}
\label{sec:interf_reference_system_tiling}
\autoref{tab:int_ref_system_slicing} reports the reference system, the slicing, and grid for interferometric products.

\begin{longtblr}[
    label=tab:int_ref_system_slicing,
    entry = none,
    caption = {Reference system and Slicing for interferometric products},
    ]{
        width=\linewidth,
        colspec = {Q[l,20]Q[l,35]Q[l,30]Q[l,15]},
        %vline{2-5} = {-}{},
        %vline{5-5} = {-}{},
        hlines,
        hline{1-Z} = {-}{0.05em},
        rowhead=2
    }
    \textbf{Product} & \textbf{Reference System} & \textbf{Slicing} & \textbf{Spacing}\\
    \textbf{\acs{int}+\acs{cohe}} (default) & Radar: azimuth/slant-range & Burst & ~50 m \\
    \textbf{\acs{int}+\acs{cohe}} (optional) & Radar: azimuth/slant-range & Burst & Full radar resolution \\
    \textbf{\acs{unw}+\acs{cohe}} (default) & Radar: azimuth/slant-range & Burst & ~50 m \\
    \textbf{\acs{unw}+\acs{cohe}} (optional) & Radar: azimuth/slant-range & Burst & Full radar resolution \\ 
    \textbf{\acs{gint}+\acs{gcohe}} & Cartographic: UTM/UPS & Burst-projected & 50 m \\
    \textbf{\acs{gunw}+\acs{gcohe}} & Cartographic: UTM/UPS & Burst-projected & 50 m \\
\end{longtblr}

\subsection{Product Requirements}
\subsubsection{Accuracy}
There is no standard praxis for defining requirements for phase products. The most natural approach would be to derive them from user requirements at the application level.

Requirements on product accuracy should ideally be specified in terms of bias and variance of the estimated phase or mean velocity. This is not trivial, as it depends on the target's correlation characteristics, and user requirements typically do not provide this level of detail.

Therefore, for this set of products, it is not straightforward to define numerical performance requirements. However, several key elements can be identified that critically influence the overall quality of the derived interferometric outputs.

The expected performance of interferometric products—specifically, the interferogram and coherence map—is primarily governed by three main factors: the accuracy of the co-registration process, the quality of the orbital data and Digital Elevation Model (\acs{dem}) used for phase flattening, and the effectiveness of atmospheric correction strategies in mitigating residual propagation delays.

First, achieving sub-pixel co-registration accuracy is essential to preserve phase coherence between image pairs. As a general guideline, a co-registration error lower than 10\% of a resolution cell is required to avoid significant decorrelation. Misalignments between images lead to spatial decorrelation which directly degrades the quality of interferometric measurements.

Second, the precision of the orbital state vectors and the spatial resolution and vertical accuracy of the \acs{dem} are critical for effective phase flattening. Inaccuracies in these inputs can introduce residual phase ramps or localized artifacts, reducing the fidelity of the interferometric phase and potentially compromising subsequent geophysical analyses.

Third, atmospheric disturbances—such as ionospheric phase delays in L-band systems or tropospheric variability in C-band—can introduce spatially variable phase distortions. The ability to estimate and correct these residuals, whether through external atmospheric models or empirical, data-driven approaches, is essential for preserving the spatial consistency of both the interferometric phase and the derived coherence.

Collectively, these factors determine the quality and usability of interferometric products and must be evaluated with respect to the sensor characteristics, acquisition geometry, and the requirements of the specific application.

When considering the unwrapped phase product derived from the interferogram, additional challenges arise, making it even more difficult to define numerical quality indicators for this type of \acs{ard}. The quality and reliability of the unwrapped phase are, in addition to the previously described elements, also strongly influenced by the spatial and temporal characteristics of the ground deformation being observed.

In particular, deformation signals with abrupt spatial gradients (e.g., localized subsidence, fault ruptures) or rapid temporal evolution (e.g., coseismic events or landslides) can significantly affect the robustness of phase unwrapping algorithms. Such conditions often lead to areas with high phase gradients, where the unwrapping process becomes more susceptible to errors, resulting in  discontinuities or unwrapping ambiguities.

Common unwrapping algorithms, such as the Minimum Cost Flow (MCF) and the Statistical-Cost Network-Flow Algorithm for Phase Unwrapping (SNAPHU), rely on assumptions of spatial smoothness and phase continuity. When these assumptions are violated due to complex deformation patterns or low-coherence areas, the algorithms may fail to recover a consistent and continuous deformation signal.

For all the geocoded \acs{ard}s—namely, \acs{gint}, \acs{gcohe}, and \acs{gunw}—which correspond respectively to the interferogram phase (\acs{int}), coherence (\acs{cohe}), and unwrapped phase (\acs{unw}) \acs{ard} products in radar geometry, the quality levels shall be equivalent to those of their radar-coordinate counterparts. This equivalence shall be ensured by applying a geocoding algorithm that preserves the intrinsic phase and amplitude information, reprojecting the data onto the ground reference grid without introducing radiometric or geometric distortions.

As with the previously discussed quality elements, it is also difficult to define a strict numerical requirement for geolocation accuracy, as it depends on several external factors. In particular, the requirements applicable to geolocation accuracy are directly dependent on the quality of the orbital data and the \acs{dem} used to project ground coordinates onto the radar geometry. The accuracy achievable by the geocoding procedure is linked to the precision of this mapping process, which depends on the spatial and vertical accuracy of the \acs{dem} and the precision of the orbit information employed.

\begin{itemize}
\item phase continuity between adjacent bursts
\item fringes in correspondence of flat terrain
\end{itemize}

\subsubsection{Timeliness}
In terms of timeliness, the \acs{int} product generation is closely related to the generation of \acs{cslc} products, considering all calibration operations performed at \acs{cslc} level. In fact, the successive phases of filtering, interferogram and coherence computation, and even the most demanding process of phase unwrapping, do not significantly impact the product delivery time.

On the contrary, if corrections are not applied during \acs{cslc} product generation, additional time must be considered for the retrieval of atmospheric models or the application of data-driven techniques for ionospheric phase screen removal to provide the \acs{int} product.

%\hl{(required? look at applications? where? ROSE-L MRD? NISAR?)
%ROSE-L MRD mentions now 1h latency in case of emergencies. This requirement might be very difficult to achieve.}

\subsection{Usability, Integrability and Interoperability}
\label{sec:usability_integrability_InSAR}
The characteristics of the interferometric products in terms of usability, integrability and interoperability are in the following table.

\begin{longtblr}[
  caption = {Usability, Integrability, and Interoperability of Interferometric Products},
  label = {tab:insar-usability-integrability},
]{
  colspec = {Q[2cm] Q[12cm]},
  rowhead = 1,
  width = \linewidth
}
\hline
\textbf{Aspect} & \textbf{Description} \\
\hline
Usability &
\begin{itemize}
  \item Ideal for detecting ground motion such as fast landslides, subsidence, earthquakes, glacier deformation.
  \item Sensitive to vertical and east-west displacements.
  \item Applicable at multiple scales: Suitable for both local and regional/global deformation monitoring.
  \item \acs{cohe} and \acs{gcohe} products support change detection, feature extraction, and thematic classification.
\end{itemize} 
\\
Integrability &
\begin{itemize}
  \item Interdisciplinary use: Applicable in geoscience, civil engineering, urban planning, agriculture, and disaster risk management.
  \item Workflow extensibility: Seamlessly embedded in advanced processing chains for ground deformation estimation.
\end{itemize} 
\\
Interoperability &
\begin{itemize}
  \item Geospatial alignment: Aligned to shared grids and geographic reference systems, allowing overlay with optical, DEM, and other vector layers.
  \item Cross-mission consistency: Enables joint use of products from multiple missions (e.g., Sentinel-1 and ROSE-L) using common radar grids.
\end{itemize} 
\\
\hline
\end{longtblr}